\documentclass[11pt,a4paper]{article}
\usepackage[utf8]{inputenc}
\usepackage[T1]{fontenc}
\usepackage[french]{babel}
\usepackage{amsmath,amssymb,amsfonts,amsthm,mathrsfs,mathtools}
\usepackage[top=2cm,bottom=2cm,left=2cm,right=2cm]{geometry}
\usepackage[dvipsnames]{xcolor}
\usepackage{tcolorbox}
\tcbuselibrary{skins,breakable}
\usepackage{tikz}
\usetikzlibrary{calc,shapes,arrows.meta}
\usepackage{pgfplots}
\pgfplotsset{compat=1.18}
\usepackage{enumitem}
\usepackage{fancyhdr}
\usepackage{lastpage}
\usepackage{hyperref}

% ----------------------------------------------------
% Couleurs officielles du projet (Charte LLG / MPSI)
% ----------------------------------------------------
\definecolor{llgblue}{RGB}{20, 50, 115}
\definecolor{llglight}{RGB}{246, 249, 255}

\definecolor{emeraldframe}{RGB}{18, 110, 75}
\definecolor{emeraldback}{RGB}{242, 250, 246}
\definecolor{emeraldtitle}{RGB}{12, 80, 52}

\definecolor{prepaframe}{RGB}{30, 70, 135}
\definecolor{prepaback}{RGB}{246, 249, 255}
\definecolor{prepatitle}{RGB}{20, 50, 105}

\definecolor{alertframe}{RGB}{185, 75, 20}
\definecolor{alertback}{RGB}{255, 249, 245}
\definecolor{alerttitle}{RGB}{145, 55, 12}

\definecolor{bonusframe}{RGB}{120, 40, 140}
\definecolor{bonusback}{RGB}{252, 247, 255}
\definecolor{bonustitle}{RGB}{90, 25, 110}

% Configuration des hyperliens
\hypersetup{
    colorlinks=true,
    linkcolor=llgblue,
    urlcolor=llgblue,
    citecolor=llgblue
}

% En-têtes et pieds de page
\pagestyle{fancy}
\fancyhf{}
\renewcommand{\headrulewidth}{0.5pt}
\renewcommand{\footrulewidth}{0.4pt}
\fancyhead[L]{\small\textbf{Lycée Louis-le-Grand} -- 1\textsuperscript{ère} Spécialité Mathématiques (1E3)}
\fancyhead[C]{\small\textbf{DST 4 : Dérivation \& Applications}}
\fancyhead[R]{\small 27 Janvier 2022}
\fancyfoot[L]{\footnotesize LLG 2021/2022 -- Devoir Sur Table n°4}
\fancyfoot[C]{\footnotesize Énoncé \& Corrigé Intégral Rigoureux}
\fancyfoot[R]{\footnotesize Page \thepage\ / \pageref{LastPage}}

% ----------------------------------------------------
% Boîtes tcolorbox personnalisées
% ----------------------------------------------------
\newtcolorbox{correctionbox}[1][]{
  colback=emeraldback,
  colframe=emeraldframe,
  fonttitle=\bfseries\color{white},
  coltitle=white,
  colbacktitle=emeraldtitle,
  title={Corrigé officiel détaillé et barème indicatif #1},
  arc=2mm,
  boxrule=0.8pt,
  titlerule=0.5pt,
  breakable,
  before skip=0.4cm,
  after skip=0.5cm
}

\newtcolorbox{remarquebox}[1][]{
  colback=prepaback,
  colframe=prepaframe,
  fonttitle=\bfseries\small\color{white},
  coltitle=white,
  colbacktitle=prepatitle,
  title={Prolongement \& Regard Prépa / MPSI #1},
  arc=1.5mm,
  boxrule=0.6pt,
  titlerule=0.4pt,
  breakable,
  before skip=0.25cm,
  after skip=0.35cm
}

\newtcolorbox{alertebox}[1][]{
  colback=alertback,
  colframe=alertframe,
  fonttitle=\bfseries\small\color{white},
  coltitle=white,
  colbacktitle=alerttitle,
  title={Point de vigilance méthodologique \& Pièges fréquents #1},
  arc=1.5mm,
  boxrule=0.6pt,
  titlerule=0.4pt,
  breakable,
  before skip=0.25cm,
  after skip=0.35cm
}

\newtcolorbox{bonusbox}[1][]{
  colback=bonusback,
  colframe=bonusframe,
  fonttitle=\bfseries\color{white},
  coltitle=white,
  colbacktitle=bonustitle,
  title={Question Bonus (+2 points) : Résolution multicritère #1},
  arc=2mm,
  boxrule=0.8pt,
  titlerule=0.5pt,
  breakable,
  before skip=0.3cm,
  after skip=0.4cm
}

% Commandes mathématiques
\newcommand{\N}{\mathbb{N}}
\newcommand{\Z}{\mathbb{Z}}
\newcommand{\R}{\mathbb{R}}
\newcommand{\C}{\mathbb{C}}
\newcommand{\dd}{\mathrm{d}}

\begin{document}

% ====================================================
% EN-TÊTE OFFICIEL DU DEVOIR
% ====================================================
\begin{center}
    {\LARGE \bfseries Lycée Louis-le-Grand -- Paris}\\[0.2cm]
    {\large \textbf{Classe de Première Spécialité Mathématiques -- Classe 1E3}}\\[0.15cm]
    {\normalsize Année scolaire 2021--2022}\\[0.35cm]
    \hrule height 1.2pt
    \vspace{0.25cm}
    {\LARGE \bfseries DEVOIR SUR TABLE N°4 : DÉRIVATION}\\[0.15cm]
    {\large \textbf{Étude locale, Équations de tangentes, Dérivabilité \& Faisceaux de courbes}}\\[0.2cm]
    \hrule height 0.6pt
    \vspace{0.25cm}
    {\small \textbf{Date :} Jeudi 27 Janvier 2022 \quad $\vert$ \quad \textbf{Durée :} 2 heures \quad $\vert$ \quad \textbf{Calculatrice interdite}}\\[0.1cm]
    {\small \textbf{Barème global :} 20 points \quad $\vert$ \quad \textbf{Bonus :} +2 points}
\end{center}

\vspace{0.3cm}

\begin{tcolorbox}[colback=llglight,colframe=llgblue,arc=1.5mm,boxrule=0.7pt]
\textbf{Consignes officielles d'évaluation :}
\begin{itemize}[leftmargin=*,itemsep=2pt]
    \item La qualité de la rédaction, la clarté des raisonnements et la précision des justifications mathématiques jouent un rôle déterminant dans l'évaluation.
    \item Tout résultat énoncé sans justification rigoureuse (théorèmes généraux de dérivation, formes indéterminées levées par le taux d'accroissement, conditions nécessaires et suffisantes de tangence) sera considéré comme nul.
    \item Les expressions doivent être simplifiées, factorisées et présentées sous leur forme la plus explicite possible.
    \item L'usage de la calculatrice, des téléphones et de tout document extérieur est formellement interdit.
\end{itemize}
\end{tcolorbox}

\vspace{0.4cm}

% ####################################################
% PARTIE I : ÉNONCÉS INTÉGRAUX (SUJET DU DST)
% ####################################################
\begin{center}
    \rule{0.35\textwidth}{0.8pt}\\[0.15cm]
    {\Large \bfseries PREMIÈRE PARTIE : ÉNONCÉ DU SUJET}\\[0.1cm]
    \rule{0.35\textwidth}{0.8pt}
\end{center}

\vspace{0.2cm}

% ----------------------------------------------------
% EXERCICE 1 (ÉNONCÉ)
% ----------------------------------------------------
\subsection*{Exercice 1 \hfill \normalfont\normalsize (5 points)}
\textbf{Identification d'un polynôme de degré 3 par conditions d'incidence et de contact}

\medskip
Soit une fonction $f$ polynomiale de degré 3, définie sur $\R$. On se place dans le plan muni d'un repère orthogonal et on note $\mathcal{C}_f$ sa courbe représentative.

On suppose que la courbe $\mathcal{C}_f$ vérifie l'ensemble des trois propriétés géométriques suivantes :
\begin{enumerate}[label=\textbf{\arabic*.}]
    \item $\mathcal{C}_f$ coupe l'axe des ordonnées au point d'ordonnée $1$.
    \item $\mathcal{C}_f$ passe par le point $A(-1\,;\,3)$ et possède en ce point une tangente de coefficient directeur égal à $-3$.
    \item $\mathcal{C}_f$ possède une tangente horizontale au point d'abscisse $0$.
\end{enumerate}

\medskip
Déterminer l'expression explicite de la fonction $f$.

\vspace{0.4cm}

% ----------------------------------------------------
% EXERCICE 2 (ÉNONCÉ)
% ----------------------------------------------------
\subsection*{Exercice 2 \hfill \normalfont\normalsize (4{,}5 points)}
\textbf{Bitangente à une quartique et contact d'ordre supérieur}

\medskip
Soit $f$ la fonction définie sur $\R$ par :
\[
f(x) = x^4 - 2x^3 - 3x^2.
\]
On note $\mathcal{C}_f$ sa courbe représentative dans un repère du plan.

\begin{enumerate}
    \item Calculer la fonction dérivée $f'(x)$ pour tout réel $x$. \hfill \textbf{(0{,}75 point)}
    \item On considère la droite $T_{-1}$, tangente à la courbe $\mathcal{C}_f$ au point d'abscisse $-1$. Montrer que l'équation réduite de $T_{-1}$ est :
    \[
    y = -4x - 4.
    \] \hfill \textbf{(1 point)}
    \item Démontrer que la droite $T_{-1}$ est également tangente à la courbe $\mathcal{C}_f$ en un second point distinct, dont on déterminera avec précision les coordonnées. \hfill \textbf{(2{,}75 points)}
\end{enumerate}

\vspace{0.4cm}

% ----------------------------------------------------
% EXERCICE 3 (ÉNONCÉ)
% ----------------------------------------------------
\subsection*{Exercice 3 \hfill \normalfont\normalsize (4{,}5 points)}
\textbf{Dérivabilité sous radicaux et étude fine aux bornes du domaine}

\medskip
On considère la fonction $f$ définie par l'expression :
\[
f(x) = (x^2 - 1)\sqrt{x - 1}.
\]
\begin{enumerate}
    \item Déterminer le domaine de définition de la fonction $f$, noté $\mathcal{D}_f$. \hfill \textbf{(0{,}5 point)}
    \item Justifier avec rigueur que $f$ est dérivable sur l'intervalle ouvert $]1\,;\,+\infty[$. \hfill \textbf{(0{,}75 point)}
    \item Donner l'expression factorisée et simplifiée de la fonction dérivée $f'(x)$ pour tout $x \in ]1\,;\,+\infty[$. \hfill \textbf{(1{,}25 point)}
    \item Étudier la dérivabilité de $f$ en $x = 1$. Interpréter géométriquement le résultat obtenu pour la courbe représentative $\mathcal{C}_f$. \hfill \textbf{(1{,}25 point)}
    \item Justifier que la formule $f'(x) = \dfrac{1}{2}(5x + 1)\sqrt{x - 1}$ est en réalité valable sur l'intervalle fermé $[1\,;\,+\infty[$. \hfill \textbf{(0{,}75 point)}
\end{enumerate}

\vspace{0.4cm}

% ----------------------------------------------------
% EXERCICE 4 (ÉNONCÉ)
% ----------------------------------------------------
\subsection*{Exercice 4 \hfill \normalfont\normalsize (6 points)}
\textbf{Faisceau de courbes rationnelles et concourance remarquable des tangentes}

\medskip
Pour tout nombre réel $\lambda \in \R$, on considère la fonction $f_\lambda$ définie sur $\R$ par :
\[
f_\lambda(x) = \frac{x + \lambda}{x^2 + 1}.
\]
On note $\mathcal{C}_\lambda$ la courbe représentative de la fonction $f_\lambda$ dans un repère orthonormal.

\begin{enumerate}
    \item Démontrer que pour tout réel $a \in \R$, la tangente $T_{a,\lambda}$ à la courbe $\mathcal{C}_\lambda$ au point d'abscisse $a$ admet pour équation cartésienne réduite :
    \[
    T_{a,\lambda} : y = \frac{-a^2 - 2a\lambda + 1}{(a^2 + 1)^2} \, x + \frac{2a^3 + 3a^2\lambda + \lambda}{(a^2 + 1)^2}.
    \] \hfill \textbf{(2{,}5 points)}
    \item Déterminer la valeur du réel $a$ pour laquelle les tangentes aux points d'abscisse $a$ aux courbes $\mathcal{C}_\lambda$ sont mutuellement parallèles, quelle que soit la valeur attribuée au paramètre réel $\lambda$. \hfill \textbf{(1 point)}
    \item On suppose dans cette question que le réel $a$ est non nul ($a \neq 0$).
    \begin{enumerate}
        \item Démontrer que les tangentes aux points d'abscisse $a$ aux courbes $\mathcal{C}_\lambda$ sont toutes concourantes en un même point $P_a$ dont on exprimera l'abscisse uniquement en fonction de $a$. \hfill \textbf{(1{,}5 point)}
        \item Déterminer les coordonnées exactes du point de concours des tangentes lorsque $a = 1$. \hfill \textbf{(1 point)}
    \end{enumerate}
\end{enumerate}

\vspace{0.4cm}

% ----------------------------------------------------
% EXERCICE BONUS (ÉNONCÉ)
% ----------------------------------------------------
\subsection*{Exercice Bonus \hfill \normalfont\normalsize (+2 points)}
\textbf{Taux d'accroissement bivarié et signification géométrique}

\medskip
Soit $f$ une fonction définie et dérivable sur $\R$, et soit $a$ un réel fixé.

Déterminer la valeur de la limite suivante :
\[
\lim_{x \to a} \frac{x f(a) - a f(x)}{x - a}.
\]
Préciser le sens géométrique de la quantité obtenue par rapport à la courbe $\mathcal{C}_f$.

\newpage

% ####################################################
% PARTIE II : CORRIGÉS DÉTAILLÉS ET APPROFONDISSEMENTS
% ####################################################
\begin{center}
    \rule{0.35\textwidth}{0.8pt}\\[0.15cm]
    {\Large \bfseries SECONDE PARTIE : CORRIGÉ OFFICIEL \& REGARD PRÉPA}\\[0.1cm]
    \rule{0.35\textwidth}{0.8pt}
\end{center}

\vspace{0.3cm}

% ====================================================
% CORRIGÉ EXERCICE 1
% ====================================================
\subsection*{Exercice 1 : Identification d'un polynôme de degré 3 \hfill \normalfont\normalsize (5 points)}

\begin{correctionbox}
\textbf{1. Modélisation polynomiale et dérivée générale :}

Puisque $f$ est une fonction polynomiale de degré 3, il existe un quadruplet unique de réels $(a, b, c, d) \in \R^4$ avec \textbf{$a \neq 0$} tel que pour tout réel $x$ :
\[
f(x) = ax^3 + bx^2 + cx + d.
\]
En tant que fonction polynomiale, $f$ est indéfiniment dérivable sur $\R$. Sa fonction dérivée $f'$ est donnée pour tout $x \in \R$ par :
\[
f'(x) = 3ax^2 + 2bx + c.
\]
\textbf{[0{,}5 point]}

\medskip
\textbf{2. Exploitation des conditions en $x = 0$ :}
\begin{itemize}[leftmargin=*]
    \item \textbf{Propriété 1 :} La courbe $\mathcal{C}_f$ coupe l'axe des ordonnées au point d'ordonnée $1$.
    Le point de contact a pour abscisse $0$ et pour ordonnée $1$, ce qui se traduit par $f(0) = 1$. Or :
    \[
    f(0) = a(0)^3 + b(0)^2 + c(0) + d = d.
    \]
    On en déduit immédiatement :
    \[
    \mathbf{d = 1}.
    \] \hfill \textbf{[0{,}5 point]}
    \item \textbf{Propriété 3 :} La courbe $\mathcal{C}_f$ possède une tangente horizontale au point d'abscisse $0$.
    Le coefficient directeur de la tangente en $0$ est égal à $f'(0)$. La tangente étant horizontale, sa pente est nulle, soit $f'(0) = 0$. Or :
    \[
    f'(0) = 3a(0)^2 + 2b(0) + c = c.
    \]
    On en déduit immédiatement :
    \[
    \mathbf{c = 0}.
    \] \hfill \textbf{[0{,}5 point]}
\end{itemize}
L'expression de la fonction et de sa dérivée se réduisent donc à :
\[
f(x) = ax^3 + bx^2 + 1 \quad \text{et} \quad f'(x) = 3ax^2 + 2bx.
\]

\medskip
\textbf{3. Exploitation des conditions au point $A(-1\,;\,3)$ :}
\begin{itemize}[leftmargin=*]
    \item \textbf{Appartenance du point $A$ :} Puisque $A(-1\,;\,3) \in \mathcal{C}_f$, on a $f(-1) = 3$.
    \[
    f(-1) = a(-1)^3 + b(-1)^2 + 1 = -a + b + 1.
    \]
    Par conséquent :
    \[
    -a + b + 1 = 3 \iff \mathbf{-a + b = 2} \quad (L_1).
    \] \hfill \textbf{[0{,}75 point]}
    \item \textbf{Pente de la tangente en $A$ :} Le coefficient directeur de la tangente en $A$ est égal à $f'(-1)$. L'énoncé indique que cette pente vaut $-3$, d'où $f'(-1) = -3$.
    \[
    f'(-1) = 3a(-1)^2 + 2b(-1) = 3a - 2b.
    \]
    Par conséquent :
    \[
    \mathbf{3a - 2b = -3} \quad (L_2).
    \] \hfill \textbf{[0{,}75 point]}
\end{itemize}

\medskip
\textbf{4. Résolution du système linéaire :}

Les coefficients $a$ et $b$ vérifient le système linéaire $(S)$ de deux équations à deux inconnues :
\[
(S) : \begin{cases}
-a + b = 2 & (L_1) \\
3a - 2b = -3 & (L_2)
\end{cases}
\]
Multiplions l'équation $(L_1)$ par $2$ et ajoutons-la à $(L_2)$ :
\[
2(L_1) + (L_2) \implies 2(-a + b) + (3a - 2b) = 2(2) + (-3) \iff a = 1.
\]
En reportant $a = 1$ dans l'équation $(L_1)$ :
\[
-(1) + b = 2 \iff b = 2 + 1 \iff b = 3.
\]
On constate que $a = 1 \neq 0$, ce qui confirme que le degré de $f$ est bien égal à 3. \hfill \textbf{[1{,}5 point]}

\medskip
\textbf{5. Synthèse et conclusion :}

L'unique fonction polynomiale de degré 3 satisfaisant aux conditions de l'énoncé est :
\[
\mathbf{f(x) = x^3 + 3x^2 + 1}.
\] \hfill \textbf{[0{,}5 point]}
\end{correctionbox}

\begin{alertebox}[Condition de degré et piège de la tangente]
\begin{enumerate}
    \item \textbf{Ne pas omettre $a \neq 0$ :} Si la résolution du système avait conduit à $a = 0$, la fonction obtenue aurait été du second degré, ce qui aurait constitué une contradiction avec l'énoncé. Il est impératif de valider cette hypothèse en fin de calcul.
    \item \textbf{Tangente au point $A$ :} L'équation de la tangente en $A(-1\,;\,3)$ s'écrit $y = f'(-1)(x - (-1)) + f(-1) = -3(x + 1) + 3 = -3x$. Elle passe donc par l'origine du repère $O(0\,;\,0)$.
\end{enumerate}
\end{alertebox}

\begin{center}
\begin{tikzpicture}[scale=1.1, cap=round, >=Latex]
    % Grille et repère
    \draw[very thin, gray!25] (-3.5,-1) grid (2,6.5);
    \draw[->, thick] (-3.5,0) -- (2.3,0) node[right] {\small $x$};
    \draw[->, thick] (0,-1) -- (0,6.8) node[above] {\small $y$};
    \node[below left] at (0,0) {\small $O$};
    \draw[thick] (1,0.08) -- (1,-0.08) node[below] {\small $1$};
    \draw[thick] (0.08,1) -- (-0.08,1) node[left] {\small $1$};

    % Courbe f(x) = x^3 + 3x^2 + 1 sur [-3.2, 1.15]
    \draw[domain=-3.2:1.15, samples=120, very thick, llgblue] 
        plot (\x, {\x*\x*\x + 3*\x*\x + 1}) 
        node[above right] {\small $\mathcal{C}_f : y = x^3 + 3x^2 + 1$};

    % Tangente horizontale en (0, 1)
    \draw[<->, thick, ForestGreen] (-1.2, 1) -- (1.2, 1) 
        node[right] {\footnotesize Tangente horizontale : $y = 1$};
    \filldraw[ForestGreen] (0, 1) circle (2.2pt) node[below right=2pt] {\footnotesize $(0, 1)$};

    % Tangente en A(-1, 3) : pente -3, équation y = -3x
    \draw[<->, thick, alertframe] (-1.8, 5.4) -- (-0.2, 0.6) 
        node[below right] {\footnotesize $T_A : y = -3x$};
    \filldraw[alertframe] (-1, 3) circle (2.2pt) node[left=2pt] {\footnotesize $A(-1, 3)$};

    % Extremum local en (-2, 5)
    \draw[<->, dashed, gray!70] (-2.8, 5) -- (-1.2, 5) 
        node[above] {\footnotesize Max local $(-2, 5)$};
    \filldraw[gray!70] (-2, 5) circle (1.8pt);

    % Lignes de rappel
    \draw[dotted, thick, alertframe] (-1, 0) -- (-1, 3);
    \node[below] at (-1, 0) {\small $-1$};
    \draw[dotted, thick, alertframe] (0, 3) -- (-1, 3);
    \node[right] at (0, 3) {\small $3$};
\end{tikzpicture}
\end{center}

\begin{remarquebox}[Interpolation d'Hermite et dimension finie]
D'un point de vue algébrique supérieur, le problème posé est un cas particulier d'\textbf{interpolation d'Hermite}.
L'espace vectoriel $\R_3[X]$ des polynômes réels de degré inférieur ou égal à $3$ est de dimension $4$. L'application linéaire :
\[
\Phi : \R_3[X] \longrightarrow \R^4, \quad P \longmapsto \bigl(P(0), \, P'(0), \, P(-1), \, P'(-1)\bigr)
\]
a pour matrice dans la base canonique $(1, X, X^2, X^3)$ :
\[
M = \begin{pmatrix}
1 & 0 & 0 & 0 \\
0 & 1 & 0 & 0 \\
1 & -1 & 1 & -1 \\
0 & 1 & -2 & 3
\end{pmatrix}.
\]
Le déterminant de cette matrice triangulaire par blocs se calcule immédiatement : $\det(M) = 1 \times 1 \times \begin{vmatrix} 1 & -1 \\ -2 & 3 \end{vmatrix} = 3 - 2 = 1 \neq 0$.
L'application $\Phi$ est donc un \textbf{isomorphisme linéaire}. Par conséquent, pour n'importe quel quadruplet de valeurs imposées $(\alpha, \beta, \gamma, \delta) \in \R^4$, il existe un et un seul polynôme $P \in \R_3[X]$ réalisant ces conditions de position et de tangence.
\end{remarquebox}

\vspace{0.4cm}

% ====================================================
% CORRIGÉ EXERCICE 2
% ====================================================
\subsection*{Exercice 2 : Bitangente à une quartique \hfill \normalfont\normalsize (4{,}5 points)}

\begin{correctionbox}
\textbf{1. Calcul de la fonction dérivée $f'(x)$ :}

La fonction $f : x \mapsto x^4 - 2x^3 - 3x^2$ est une fonction polynomiale, donc elle est dérivable sur $\R$. Pour tout réel $x$ :
\[
f'(x) = 4x^3 - 2 \times (3x^2) - 3 \times (2x) = 4x^3 - 6x^2 - 6x = 2x(2x^2 - 3x - 3).
\] \hfill \textbf{[0{,}75 point]}

\medskip
\textbf{2. Équation de la tangente $T_{-1}$ au point d'abscisse $-1$ :}

Calculons l'image et le nombre dérivé en $-1$ :
\[
f(-1) = (-1)^4 - 2(-1)^3 - 3(-1)^2 = 1 - 2(-1) - 3(1) = 1 + 2 - 3 = 0.
\]
\[
f'(-1) = 4(-1)^3 - 6(-1)^2 - 6(-1) = -4 - 6 + 6 = -4.
\]
La tangente $T_{-1}$ à la courbe $\mathcal{C}_f$ au point $(-1\,;\,0)$ a pour équation :
\[
y = f'(-1)\bigl(x - (-1)\bigr) + f(-1) \iff y = -4(x + 1) + 0 \iff \mathbf{y = -4x - 4}.
\] \hfill \textbf{[1 point]}

\medskip
\textbf{3. Détermination du second point de contact (bitangence) :}

\medskip
\underline{\textit{Méthode 1 : Approche algébrique par factorisation euclidienne (Contact d'ordre 2)}}

Soit la droite d'équation $y = -4x - 4$. Étudions la différence $D(x) = f(x) - (-4x - 4)$ pour tout $x \in \R$ :
\[
D(x) = x^4 - 2x^3 - 3x^2 + 4x + 4.
\]
Puisque la droite est tangente à $\mathcal{C}_f$ en $x = -1$, le réel $-1$ est une racine au moins double du polynôme $D(x)$, ce qui garantit que $D(x)$ est factorisable par $(x + 1)^2 = x^2 + 2x + 1$.
Effectuons la division euclidienne de $D(x)$ par $x^2 + 2x + 1$ :
\[
\begin{array}{r|l}
x^4 - 2x^3 - 3x^2 + 4x + 4 & x^2 + 2x + 1 \\
\cline{2-2}
-(x^4 + 2x^3 + x^2)\phantom{+4x+4} & x^2 - 4x + 4 \\
\cline{1-1}
-4x^3 - 4x^2 + 4x + 4 & \\
-(-4x^3 - 8x^2 - 4x)\phantom{+4} & \\
\cline{1-1}
4x^2 + 8x + 4 & \\
-(4x^2 + 8x + 4) & \\
\cline{1-1}
0 & 
\end{array}
\]
On observe que le quotient est $x^2 - 4x + 4 = (x - 2)^2$. Ainsi, pour tout $x \in \R$ :
\[
D(x) = f(x) - (-4x - 4) = (x + 1)^2 (x - 2)^2.
\]
Cette factorisation prouve que :
\begin{itemize}[leftmargin=*]
    \item $D(2) = 0$, donc $f(2) = -4(2) - 4 = -12$. Le point $B(2\,;\,-12)$ est commun à $\mathcal{C}_f$ et à la droite $T_{-1}$.
    \item $x = 2$ est une racine de multiplicité 2 du polynôme différence $D(x)$. Or $D'(x) = f'(x) - (-4) = f'(x) + 4$. Comme $2$ est racine double de $D$, on a $D'(2) = 0$, ce qui équivaut à $f'(2) = -4$.
\end{itemize}
La droite $T_{-1}$ est donc tangente à la courbe $\mathcal{C}_f$ au point d'abscisse $2$, dont les coordonnées sont :
\[
\mathbf{B(2\,;\,-12)}.
\]

\medskip
\underline{\textit{Méthode 2 : Identification analytique complète des tangentes}}

Soit $a \in \R$. La tangente $T_a$ à $\mathcal{C}_f$ au point d'abscisse $a$ admet pour équation :
\[
y = f'(a)(x - a) + f(a) = f'(a)x + \bigl(f(a) - a f'(a)\bigr).
\]
Pour que la droite $T_a$ coïncide avec la droite $T_{-1} : y = -4x - 4$, il faut et il suffit que leurs coefficients directeurs et leurs ordonnées à l'origine soient simultanément identiques :
\[
\begin{cases}
f'(a) = -4 & (\text{Condition de pente}) \\
f(a) - a f'(a) = -4 & (\text{Condition d'ordonnée à l'origine})
\end{cases}
\]
1. Équation sur le coefficient directeur :
\[
4a^3 - 6a^2 - 6a = -4 \iff 4a^3 - 6a^2 - 6a + 4 = 0 \iff 2a^3 - 3a^2 - 3a + 2 = 0.
\]
Sachant que $a = -1$ est racine, on factorise par $(a + 1)$ :
\[
2a^3 - 3a^2 - 3a + 2 = (a + 1)(2a^2 - 5a + 2) = 0.
\]
Le discriminant du trinôme $2a^2 - 5a + 2$ est $\Delta = (-5)^2 - 4(2)(2) = 25 - 16 = 9 = 3^2 > 0$.
Ses racines sont :
\[
a_1 = \frac{5 - 3}{4} = \frac{1}{2} \quad \text{et} \quad a_2 = \frac{5 + 3}{4} = 2.
\]
Les trois abscisses candidates où la tangente a pour pente $-4$ sont donc : $a \in \left\{-1\,;\, \dfrac{1}{2}\,;\, 2\right\}$.

2. Équation sur l'ordonnée à l'origine :
On calcule $p(a) = f(a) - a f'(a) = (a^4 - 2a^3 - 3a^2) - a(4a^3 - 6a^2 - 6a) = -3a^4 + 4a^3 + 3a^2$.
Testons les candidats obtenus :
\begin{itemize}[leftmargin=*]
    \item Pour $a = -1$ : $p(-1) = -3(1) + 4(-1) + 3(1) = -4$ (point initial).
    \item Pour $a = \frac{1}{2}$ : $p\left(\frac{1}{2}\right) = -3\left(\frac{1}{16}\right) + 4\left(\frac{1}{8}\right) + 3\left(\frac{1}{4}\right) = -\frac{3}{16} + \frac{1}{2} + \frac{3}{4} = \frac{17}{16} \neq -4$.
    (Cette droite est parallèle à $T_{-1}$ mais distincte).
    \item Pour $a = 2$ : $p(2) = -3(16) + 4(8) + 3(4) = -48 + 32 + 12 = -4$.
\end{itemize}
Le second point de contact recherché est donc bien $a = 2$, d'ordonnée $f(2) = 2^4 - 2(8) - 3(4) = -12$. \hfill \textbf{[2{,}75 points]}
\end{correctionbox}

\begin{alertebox}[Distinction cruciale : Parallélisme versus Tangence commune]
Une erreur classique consiste à s'arrêter à la condition de pente $f'(a) = -4$ et à en conclure hâtivement que les trois points d'abscisses $-1$, $1/2$ et $2$ partagent la même tangente.
La condition $f'(a) = -4$ assure uniquement que les tangentes sont \textbf{parallèles entre elles}. Pour qu'elles soient \textbf{confondues}, il est impératif d'égaler les ordonnées à l'origine (ou de vérifier que le point $(a, f(a))$ appartient à la droite).
\end{alertebox}

\begin{center}
\begin{tikzpicture}[x=1.7cm, y=0.42cm, cap=round, >=Latex]
    % Grille et repère
    \draw[very thin, gray!25] (-2,-14) grid[ystep=2, xstep=1] (3,6);
    \draw[->, thick] (-2,0) -- (3.2,0) node[right] {\small $x$};
    \draw[->, thick] (0,-14.5) -- (0,6.5) node[above] {\small $y$};
    \node[below left] at (0,0) {\small $O$};
    \foreach \x in {-1, 1, 2} {
        \draw[thick] (\x, 0.3) -- (\x, -0.3) node[above=2pt] {\footnotesize $\x$};
    }
    \foreach \y in {-12, -8, -4, 4} {
        \draw[thick] (0.06, \y) -- (-0.06, \y) node[left] {\footnotesize $\y$};
    }

    % Courbe quartique
    \draw[domain=-1.65:2.55, samples=150, very thick, llgblue] 
        plot (\x, {\x*\x*\x*\x - 2*\x*\x*\x - 3*\x*\x}) 
        node[above right] {\small $\mathcal{C}_f : y = x^4 - 2x^3 - 3x^2$};

    % Bitangente y = -4x - 4
    \draw[domain=-1.7:2.4, samples=20, very thick, alertframe] 
        plot (\x, {-4*\x - 4}) 
        node[pos=0.08, above left] {\small $T_{-1} : y = -4x - 4$};

    % Tangente parallèle en 1/2 : y = -4x + 17/16
    \draw[domain=-0.4:1.4, samples=20, dashed, gray!75, thick] 
        plot (\x, {-4*\x + 1.0625}) 
        node[right] {\footnotesize Tangente parallèle en $x=\frac{1}{2}$};
    \filldraw[gray!75] (0.5, -0.9375) circle (1.8pt);

    % Points de contact
    \filldraw[alertframe] (-1, 0) circle (2.2pt) node[above left=2pt] {\footnotesize $A(-1\,;\,0)$};
    \filldraw[alertframe] (2, -12) circle (2.2pt) node[below left=2pt] {\footnotesize $B(2\,;\,-12)$};

    % Lignes de rappel
    \draw[dotted, thick, alertframe] (2, 0) -- (2, -12) -- (0, -12);
\end{tikzpicture}
\end{center}

\begin{remarquebox}[Propriété globale remarquable : Position de la quartique]
La relation obtenue par factorisation euclidienne :
\[
\forall x \in \R, \quad f(x) - (-4x - 4) = (x + 1)^2 (x - 2)^2 \ge 0
\]
démontre que la courbe quartique $\mathcal{C}_f$ est \textbf{entièrement située au-dessus de sa bitangente} sur tout $\R$.
La droite $y = -4x - 4$ est une sécante d'appui globale qui « soutient » les deux lobes minimaux de la quartique sans jamais traverser la courbe.
De plus, la dérivée seconde de $f$ vaut :
\[
f''(x) = 12x^2 - 12x - 6 = 6(2x^2 - 2x - 1).
\]
Le discriminant de $2x^2 - 2x - 1$ est $\Delta = 4 + 8 = 12 > 0$. La dérivée seconde s'annule et change de signe en $x = \frac{1 \pm \sqrt{3}}{2}$, ce qui donne exactement deux \textbf{points d'inflexion}. La présence d'une région concave entre deux régions convexes est la condition topologique nécessaire à l'existence d'une bitangente réelle.
\end{remarquebox}

\vspace{0.4cm}

% ====================================================
% CORRIGÉ EXERCICE 3
% ====================================================
\subsection*{Exercice 3 : Dérivabilité sous radicaux et étude aux bornes \hfill \normalfont\normalsize (4{,}5 points)}

\begin{correctionbox}
\textbf{1. Domaine de définition de la fonction $f$ :}

La fonction racine carrée $t \mapsto \sqrt{t}$ est définie sur $[0\,;\,+\infty[$. L'expression sous le radical doit être positive ou nulle :
\[
x - 1 \ge 0 \iff x \ge 1.
\]
La fonction polynomiale $x \mapsto x^2 - 1$ étant définie sur $\R$, le domaine de définition de $f$ est l'intervalle fermé :
\[
\mathbf{\mathcal{D}_f = [1\,;\,+\infty[}.
\] \hfill \textbf{[0{,}5 point]}

\medskip
\textbf{2. Dérivabilité sur $]1\,;\,+\infty[$ :}

Justifions la dérivabilité sur l'ouvert $]1\,;\,+\infty[$ en appliquant les théorèmes généraux :
\begin{itemize}[leftmargin=*]
    \item La fonction $u : x \mapsto x^2 - 1$ est polynomiale, donc dérivable sur $\R$, en particulier sur $]1\,;\,+\infty[$.
    \item La fonction affine $w : x \mapsto x - 1$ est dérivable sur $\R$ et vérifie $w(x) = x - 1 > 0$ pour tout $x \in ]1\,;\,+\infty[$.
    \item La fonction racine carrée $g : t \mapsto \sqrt{t}$ est dérivable sur $\R^{+*} = ]0\,;\,+\infty[$.
    \item Par composition de deux fonctions dérivables, la fonction $v : x \mapsto \sqrt{x - 1} = g\bigl(w(x)\bigr)$ est dérivable sur $]1\,;\,+\infty[$.
    \item Par produit de deux fonctions dérivables, la fonction $f = u \times v$ est rigoureusement \textbf{dérivable sur $]1\,;\,+\infty[$}.
\end{itemize} \hfill \textbf{[0{,}75 point]}

\medskip
\textbf{3. Calcul et simplification de $f'(x)$ sur $]1\,;\,+\infty[$ :}

Pour tout $x \in ]1\,;\,+\infty[$, les dérivées de $u$ et $v$ sont :
\[
u'(x) = 2x \quad \text{et} \quad v'(x) = \frac{1}{2\sqrt{x - 1}}.
\]
Appliquons la règle de dérivation du produit $(uv)' = u'v + uv'$ :
\[
f'(x) = 2x\sqrt{x - 1} + (x^2 - 1) \times \frac{1}{2\sqrt{x - 1}} = \frac{2x \times 2(x - 1) + (x^2 - 1)}{2\sqrt{x - 1}} = \frac{4x(x - 1) + (x^2 - 1)}{2\sqrt{x - 1}}.
\]
Factorisons $x^2 - 1 = (x - 1)(x + 1)$ au numérateur :
\[
4x(x - 1) + (x - 1)(x + 1) = (x - 1)\bigl(4x + x + 1\bigr) = (x - 1)(5x + 1).
\]
D'où :
\[
f'(x) = \frac{(x - 1)(5x + 1)}{2\sqrt{x - 1}}.
\]
Puisque $x > 1$, on a $x - 1 > 0$, donc $\dfrac{x - 1}{\sqrt{x - 1}} = \sqrt{x - 1}$. On obtient l'expression simplifiée remarquable :
\[
\mathbf{f'(x) = \frac{1}{2}(5x + 1)\sqrt{x - 1}}.
\] \hfill \textbf{[1{,}25 point]}

\medskip
\textbf{4. Étude de la dérivabilité en $x = 1$ et interprétation géométrique :}

La fonction racine carrée n'étant pas dérivable en $0$, la dérivabilité de $f$ en $1$ ne peut pas être déduite des théorèmes d'algèbre de dérivées. On doit impérativement revenir à la définition fondamentale du nombre dérivé par le \textbf{taux d'accroissement}.
On a $f(1) = (1^2 - 1)\sqrt{1 - 1} = 0$. Pour tout $h > 0$ :
\[
\tau_1(h) = \frac{f(1 + h) - f(1)}{h} = \frac{\bigl((1 + h)^2 - 1\bigr)\sqrt{(1 + h) - 1} - 0}{h}.
\]
Développons le premier facteur :
\[
(1 + h)^2 - 1 = 1 + 2h + h^2 - 1 = 2h + h^2 = h(2 + h).
\]
En simplifiant par $h \neq 0$ :
\[
\tau_1(h) = \frac{h(2 + h)\sqrt{h}}{h} = (2 + h)\sqrt{h}.
\]
Étudions la limite quand $h$ tend vers $0$ par valeurs supérieures ($h \to 0^+$) :
\[
\lim_{h \to 0^+} (2 + h) = 2 \quad \text{et} \quad \lim_{h \to 0^+} \sqrt{h} = 0 \implies \lim_{h \to 0^+} \tau_1(h) = 2 \times 0 = 0.
\]
Cette limite est un nombre réel fini ($0 \in \R$). Par conséquent, la fonction $f$ est \textbf{dérivable à droite en $1$}, et son nombre dérivé vaut :
\[
\mathbf{f'_d(1) = 0}.
\]
\textbf{Interprétation géométrique :} La courbe représentative $\mathcal{C}_f$ admet au point de coordonnées $(1\,;\,0)$ une \textbf{demi-tangente horizontale} orientée vers la droite, portée par l'axe des abscisses d'équation $y = 0$ pour $x \ge 1$. \hfill \textbf{[1{,}25 point]}

\medskip
\textbf{5. Validité de la formule sur l'intervalle fermé $[1\,;\,+\infty[$ :}

Évaluons l'expression $\dfrac{1}{2}(5x + 1)\sqrt{x - 1}$ en $x = 1$ :
\[
\frac{1}{2}\bigl(5(1) + 1\bigr)\sqrt{1 - 1} = \frac{1}{2}(6) \times 0 = 0.
\]
Cette valeur numérique coïncide exactement avec le nombre dérivé à droite $f'_d(1) = 0$ calculé par le taux d'accroissement à la question 4.
Par conséquent, la formule :
\[
\mathbf{f'(x) = \frac{1}{2}(5x + 1)\sqrt{x - 1}}
\]
est parfaitement continue et rigoureusement valable sur l'intervalle fermé $[1\,;\,+\infty[$. \hfill \textbf{[0{,}75 point]}
\end{correctionbox}

\begin{alertebox}[Le piège de la non-dérivabilité de la racine carrée]
Une erreur fatale consisterait à écrire : « $\sqrt{x-1}$ n'est pas dérivable en $1$, donc $f$ n'est pas dérivable en $1$ ».
C'est un faux théorème : le produit d'une fonction dérivable s'annulant en un point et d'une fonction non dérivable en ce même point peut tout à fait être dérivable ! Le facteur polynomial $(x^2 - 1) = (x-1)(x+1)$ s'annule à l'ordre 1 en $1$, ce qui compense la singularité en $1/\sqrt{h}$ et amortit le taux de variation jusqu'à lui donner une limite nulle.
\end{alertebox}

\begin{center}
\begin{tikzpicture}[x=3cm, y=0.9cm, cap=round, >=Latex]
    % Grille et repère
    \draw[very thin, gray!25] (0.6,-0.5) grid[xstep=0.5, ystep=1] (3.2,7);
    \draw[->, thick] (0.6,0) -- (3.3,0) node[right] {\small $x$};
    \draw[->, thick] (0.8,-0.5) -- (0.8,7.2) node[above] {\small $y$};

    \foreach \x in {1, 2, 3} {
        \draw[thick] (\x, 0.1) -- (\x, -0.1) node[below] {\small $\x$};
    }
    \foreach \y in {1, 2, 3, 4, 5, 6} {
        \draw[thick] (0.82, \y) -- (0.78, \y) node[left] {\small $\y$};
    }

    % Courbe f(x)
    \draw[domain=1:2.65, samples=150, very thick, llgblue] 
        plot (\x, {(\x*\x - 1)*sqrt(\x - 1)}) 
        node[above left] {\small $\mathcal{C}_f : y = (x^2-1)\sqrt{x-1}$};

    % Demi-tangente horizontale en (1,0)
    \draw[->, ultra thick, ForestGreen] (1, 0) -- (1.85, 0) 
        node[above=2pt] {\footnotesize Demi-tangente horizontale : $f'_d(1) = 0$};
    \filldraw[ForestGreen] (1, 0) circle (2.2pt) node[below left=2pt] {\footnotesize $(1,0)$};

    % Point (2, 3)
    \filldraw[llgblue] (2, 3) circle (2pt) node[below right] {\footnotesize $(2, 3)$};
    \draw[dotted, thick, gray] (2, 0) -- (2, 3) -- (0.8, 3);
\end{tikzpicture}
\end{center}

\begin{remarquebox}[Théorème de prolongement de la dérivée (MPSI)]
En classe préparatoire, ce résultat s'obtient de façon élégante grâce au \textbf{Théorème de la limite de la dérivée} (corollaire du Théorème des Accroissements Finis) :
Si une fonction $f$ est continue sur $[1\,;\,b]$, dérivable sur $]1\,;\,b[$, et si la dérivée admet une limite finie $\ell \in \R$ quand $x \to 1^+$, alors $f$ est dérivable à droite en $1$ et $f'_d(1) = \ell$.
Ici, $f$ est continue sur $[1\,;\,+\infty[$ et pour tout $x > 1$ :
\[
\lim_{x \to 1^+} f'(x) = \lim_{x \to 1^+} \frac{1}{2}(5x + 1)\sqrt{x - 1} = \frac{1}{2}(6) \times 0 = 0.
\]
On conclut immédiatement que $f$ est dérivable à droite en $1$ avec $f'_d(1) = 0$, et que $f$ est de classe $\mathcal{C}^1$ sur $[1\,;\,+\infty[$.
\end{remarquebox}

\vspace{0.4cm}

% ====================================================
% CORRIGÉ EXERCICE 4
% ====================================================
\subsection*{Exercice 4 : Faisceau de courbes et concourance des tangentes \hfill \normalfont\normalsize (6 points)}

\begin{correctionbox}
\textbf{1. Démonstration de l'équation réduite de la tangente $T_{a,\lambda}$ :}

Pour tout réel $x \in \R$, $x^2 + 1 \ge 1 > 0$. Le dénominateur ne s'annule jamais sur $\R$.
La fonction $f_\lambda : x \mapsto \dfrac{x + \lambda}{x^2 + 1}$ est une fonction rationnelle à dénominateur strictement positif, elle est donc indéfiniment dérivable sur $\R$.
Posons $u(x) = x + \lambda$ et $v(x) = x^2 + 1$. Pour tout $x \in \R$ :
\[
u'(x) = 1 \quad \text{et} \quad v'(x) = 2x.
\]
En appliquant la formule de dérivation d'un quotient $\left(\dfrac{u}{v}\right)' = \dfrac{u'v - uv'}{v^2}$ :
\[
f_\lambda'(x) = \frac{1 \cdot (x^2 + 1) - (x + \lambda)(2x)}{(x^2 + 1)^2} = \frac{x^2 + 1 - 2x^2 - 2\lambda x}{(x^2 + 1)^2} = \frac{-x^2 - 2\lambda x + 1}{(x^2 + 1)^2}.
\]
Au point d'abscisse $a$, le coefficient directeur de la tangente $T_{a,\lambda}$ est donc :
\[
f_\lambda'(a) = \frac{-a^2 - 2a\lambda + 1}{(a^2 + 1)^2}.
\]
L'ordonnée du point de contact sur la courbe $\mathcal{C}_\lambda$ est :
\[
f_\lambda(a) = \frac{a + \lambda}{a^2 + 1} = \frac{(a + \lambda)(a^2 + 1)}{(a^2 + 1)^2} = \frac{a^3 + a + a^2\lambda + \lambda}{(a^2 + 1)^2}.
\]
L'équation de la tangente $T_{a,\lambda}$ s'écrit $y = f_\lambda'(a)(x - a) + f_\lambda(a) = f_\lambda'(a)x + p(a,\lambda)$ où l'ordonnée à l'origine $p(a,\lambda)$ vaut :
\[
p(a,\lambda) = f_\lambda(a) - a f_\lambda'(a).
\]
Réduisons cette expression au même dénominateur $(a^2 + 1)^2$ :
\[
p(a,\lambda) = \frac{(a^3 + a + a^2\lambda + \lambda) - a(-a^2 - 2a\lambda + 1)}{(a^2 + 1)^2}.
\]
Développons avec soin le numérateur :
\[
a^3 + a + a^2\lambda + \lambda - (-a^3 - 2a^2\lambda + a) = a^3 + a + a^2\lambda + \lambda + a^3 + 2a^2\lambda - a = 2a^3 + 3a^2\lambda + \lambda.
\]
On obtient ainsi rigoureusement l'équation demandée :
\[
\mathbf{T_{a,\lambda} : y = \frac{-a^2 - 2a\lambda + 1}{(a^2 + 1)^2} \, x + \frac{2a^3 + 3a^2\lambda + \lambda}{(a^2 + 1)^2}}.
\] \hfill \textbf{[2{,}5 points]}

\medskip
\textbf{2. Condition de parallélisme universel des tangentes :}

Deux droites non verticales sont parallèles si et seulement si elles possèdent le même coefficient directeur.
Les tangentes $T_{a,\lambda}$ aux points d'abscisse $a$ sont parallèles entre elles pour toute valeur de $\lambda$ si et seulement si leur pente $m(a,\lambda) = f_\lambda'(a)$ est indépendante de $\lambda$.
Décomposons cette pente sous forme affine par rapport au paramètre $\lambda$ :
\[
m(a,\lambda) = \frac{1 - a^2}{(a^2 + 1)^2} + \lambda \times \left( \frac{-2a}{(a^2 + 1)^2} \right).
\]
Pour que cette quantité soit constante quel que soit $\lambda \in \R$, il faut et il suffit que le coefficient devant $\lambda$ soit identiquement nul :
\[
\frac{-2a}{(a^2 + 1)^2} = 0 \iff -2a = 0 \iff \mathbf{a = 0}.
\]
\textbf{Vérification :} Pour $a = 0$, la pente vaut $\dfrac{1}{(0^2 + 1)^2} = 1$ pour tout $\lambda \in \R$, et l'équation de la tangente devient $y = x + \lambda$. Toutes ces droites ont pour coefficient directeur 1 : elles sont mutuellement parallèles et parallèles à la première bissectrice. \hfill \textbf{[1 point]}

\medskip
\textbf{3. (a) Concourance des tangentes pour $a \neq 0$ :}

Soit $a \neq 0$. Un point $M(x\,;\,y)$ appartient à la tangente $T_{a,\lambda}$ pour toute valeur de $\lambda \in \R$ si et seulement si ses coordonnées $(x, y)$ vérifient l'équation de $T_{a,\lambda}$ quel que soit le paramètre réel $\lambda$.
Puisque $(a^2 + 1)^2 \neq 0$, multiplions l'équation par $(a^2 + 1)^2$ :
\[
(a^2 + 1)^2 y = (-a^2 - 2a\lambda + 1)x + 2a^3 + 3a^2\lambda + \lambda.
\]
Transposons tous les termes dans le membre de gauche et regroupons selon les puissances de $\lambda$ :
\[
\lambda \cdot \bigl(-2ax + 3a^2 + 1\bigr) + \Bigl[(1 - a^2)x + 2a^3 - (a^2 + 1)^2 y\Bigr] = 0.
\]
Cette équation est de la forme $A(x) \lambda + B(x,y) = 0$.
Un polynôme affine en $\lambda$ est identiquement nul sur $\R$ si et seulement si son coefficient directeur et son terme constant sont simultanément nuls :
\[
\begin{cases}
-2ax + 3a^2 + 1 = 0 & (E_1) \\
(1 - a^2)x + 2a^3 - (a^2 + 1)^2 y = 0 & (E_2)
\end{cases}
\]
L'équation $(E_1)$ détermine de manière unique l'abscisse $x$, car par hypothèse $a \neq 0$ :
\[
-2ax + 3a^2 + 1 = 0 \iff 2ax = 3a^2 + 1 \iff \mathbf{x = \frac{3a^2 + 1}{2a}}.
\]
Cette valeur de l'abscisse est strictement indépendante de $\lambda$. Cela prouve que toutes les tangentes $T_{a,\lambda}$ se coupent en un point dont l'abscisse dépend uniquement de $a$. \hfill \textbf{[1{,}5 point]}

\medskip
\textbf{3. (b) Coordonnées du point de concours pour $a = 1$ :}

Pour $a = 1$, calculons l'abscisse du point de concours :
\[
x = \frac{3(1)^2 + 1}{2(1)} = \frac{4}{2} = \mathbf{2}.
\]
Pour obtenir son ordonnée $y$, injectons $a = 1$ et $x = 2$ dans l'équation $(E_2)$ :
\[
(1 - 1^2)(2) + 2(1)^3 - (1^2 + 1)^2 y = 0 \iff 0 + 2 - (2)^2 y = 0 \iff 2 - 4y = 0 \iff \mathbf{y = \frac{1}{2}}.
\]
Ainsi, pour tout $\lambda \in \R$, la tangente à $\mathcal{C}_\lambda$ au point d'abscisse $1$ passe impérativement par le point fixe :
\[
\mathbf{P\left(2\,;\,\frac{1}{2}\right)}.
\] \hfill \textbf{[1 point]}
\end{correctionbox}

\begin{alertebox}[Méthode générale pour les faisceaux de droites]
Pour démontrer qu'une famille de droites $(D_\lambda)$ dépendant d'un paramètre $\lambda$ est concourante en un point fixe :
\begin{enumerate}
    \item Réécrire l'équation cartésienne sous forme d'un polynôme en $\lambda$ : $\lambda \, P(x,y) + Q(x,y) = 0$.
    \item Annuler simultanément $P(x,y) = 0$ et $Q(x,y) = 0$ (principe d'identification des polynômes nuls).
    \item Résoudre le système obtenu pour déterminer les coordonnées du point pivot.
\end{enumerate}
\end{alertebox}

\begin{center}
\begin{tikzpicture}[scale=1.35, cap=round, >=Latex]
    % Grille et repère
    \draw[very thin, gray!25] (-2.5,-1.5) grid[step=0.5] (3.5,2.5);
    \draw[->, thick] (-2.5,0) -- (3.7,0) node[right] {\small $x$};
    \draw[->, thick] (0,-1.5) -- (0,2.7) node[above] {\small $y$};
    \node[below left] at (0,0) {\small $O$};

    \foreach \x in {-2, -1, 1, 2, 3} {
        \draw[thick] (\x, 0.06) -- (\x, -0.06) node[below] {\footnotesize $\x$};
    }
    \foreach \y in {-1, 1, 2} {
        \draw[thick] (0.06, \y) -- (-0.06, \y) node[left] {\footnotesize $\y$};
    }

    % Courbes C_\lambda
    % \lambda = -1 (violet)
    \draw[domain=-2.5:3.2, samples=100, thick, violet!80!black] 
        plot (\x, {(\x - 1)/(\x*\x + 1)}) node[left=2pt] {\footnotesize $\mathcal{C}_{-1}$};
    % \lambda = 0 (bleu)
    \draw[domain=-2.5:3.2, samples=100, thick, llgblue] 
        plot (\x, {(\x)/(\x*\x + 1)}) node[above left=2pt] {\footnotesize $\mathcal{C}_{0}$};
    % \lambda = 1 (teal)
    \draw[domain=-2.5:3.2, samples=100, thick, teal!90!black] 
        plot (\x, {(\x + 1)/(\x*\x + 1)}) node[above right=2pt] {\footnotesize $\mathcal{C}_{1}$};
    % \lambda = 2 (orange)
    \draw[domain=-2.5:3.2, samples=100, thick, orange!90!black] 
        plot (\x, {(\x + 2)/(\x*\x + 1)}) node[right=2pt] {\footnotesize $\mathcal{C}_{2}$};

    % Tangentes en a = 1
    % \lambda = -1 : y = 0.5*x - 0.5
    \draw[domain=0.5:3.0, samples=2, dashed, violet!90!black, thick] 
        plot (\x, {0.5*\x - 0.5});
    % \lambda = 0 : y = 0.5
    \draw[domain=0.5:3.0, samples=2, dashed, llgblue, thick] 
        plot (\x, {0.5});
    % \lambda = 1 : y = -0.5*x + 1.5
    \draw[domain=0.5:3.0, samples=2, dashed, teal!90!black, thick] 
        plot (\x, {-0.5*\x + 1.5});
    % \lambda = 2 : y = -x + 2.5
    \draw[domain=0.5:3.0, samples=2, dashed, orange!90!black, thick] 
        plot (\x, {-\x + 2.5});

    % Points de contact en x = 1
    \filldraw[violet!80!black] (1, 0) circle (1.8pt);
    \filldraw[llgblue] (1, 0.5) circle (1.8pt);
    \filldraw[teal!90!black] (1, 1.0) circle (1.8pt);
    \filldraw[orange!90!black] (1, 1.5) circle (1.8pt);

    % Point de concours P(2, 0.5)
    \filldraw[alertframe] (2, 0.5) circle (2.6pt) 
        node[above right=3pt] {\bfseries\small $P\left(2\,;\,\frac{1}{2}\right)$};
    \draw[dotted, thick, alertframe] (2, 0) -- (2, 0.5) -- (0, 0.5);
\end{tikzpicture}
\end{center}

\begin{remarquebox}[Lieu géométrique des points de concours $P_a$ (MPSI)]
Que devient le point de concours $P_a\bigl(x(a)\,;\,y(a)\bigr)$ lorsque le réel $a$ varie dans $\R^*$ ?
Reprenons le système d'identification $(E_1)$ et $(E_2)$ :
\[
x(a) = \frac{3a^2 + 1}{2a} = \frac{3a}{2} + \frac{1}{2a}.
\]
De l'équation $(E_2)$, on tire :
\[
y(a) = \frac{(1 - a^2)x(a) + 2a^3}{(a^2 + 1)^2} = \frac{(1 - a^2)\frac{3a^2+1}{2a} + 2a^3}{(a^2 + 1)^2} = \frac{3a^2 + 1 - 3a^4 - a^2 + 4a^4}{2a(a^2 + 1)^2} = \frac{a^4 + 2a^2 + 1}{2a(a^2 + 1)^2} = \frac{(a^2 + 1)^2}{2a(a^2 + 1)^2} = \frac{1}{2a}.
\]
On obtient la simplification spectaculaire :
\[
y(a) = \frac{1}{2a} \iff a = \frac{1}{2y}.
\]
En reportant cette valeur dans $x(a)$ :
\[
x = \frac{3}{2}\left(\frac{1}{2y}\right) + y = \frac{3}{4y} + y \iff 4xy = 3 + 4y^2 \iff \mathbf{4xy - 4y^2 = 3}.
\]
Le lieu géométrique des points de concours $P_a$ lorsque $a$ parcourt $\R^*$ est une \textbf{hyperbole} dans le plan cartésien !
\end{remarquebox}

\vspace{0.4cm}

% ====================================================
% CORRIGÉ EXERCICE BONUS
% ====================================================
\subsection*{Exercice Bonus : Taux de variation bivarié \hfill \normalfont\normalsize (+2 points)}

\begin{bonusbox}
Soit $f$ une fonction dérivable sur $\R$ et $a \in \R$. Calculons la limite :
\[
L = \lim_{x \to a} \frac{x f(a) - a f(x)}{x - a}.
\]

\textbf{Méthode 1 : Astuce de télescopage (Définition du nombre dérivé)}

Lorsque $x \to a$, le numérateur tend vers $a f(a) - a f(a) = 0$ et le dénominateur vers $0$ : il s'agit d'une forme indéterminée « $\dfrac{0}{0}$ ».
Pour faire émerger le taux d'accroissement de $f$ en $a$, introduisons le terme croisé $a f(a)$ en l'ajoutant et le retranchant au numérateur :
\[
\frac{x f(a) - a f(x)}{x - a} = \frac{x f(a) - a f(a) + a f(a) - a f(x)}{x - a}.
\]
Factorisons par $f(a)$ d'une part, et par $-a$ d'autre part :
\[
\frac{x f(a) - a f(x)}{x - a} = \frac{(x - a)f(a) - a\bigl(f(x) - f(a)\bigr)}{x - a}.
\]
Puisque $x \neq a$, séparons la fraction en deux termes :
\[
\frac{x f(a) - a f(x)}{x - a} = \frac{(x - a)f(a)}{x - a} - a \times \frac{f(x) - f(a)}{x - a} = f(a) - a \times \frac{f(x) - f(a)}{x - a}.
\]
Par hypothèse, la fonction $f$ est dérivable en $a$. Par définition :
\[
\lim_{x \to a} \frac{f(x) - f(a)}{x - a} = f'(a).
\]
Par opérations sur les limites (somme et produit par la constante $-a$) :
\[
\mathbf{\lim_{x \to a} \frac{x f(a) - a f(x)}{x - a} = f(a) - a f'(a)}.
\] \hfill \textbf{[1 point]}

\medskip
\textbf{Interprétation géométrique fondamentale :}

Rappelons l'équation de la tangente $T_a$ à la courbe $\mathcal{C}_f$ au point d'abscisse $a$ :
\[
y = f'(a)(x - a) + f(a) = f'(a)x + \bigl(f(a) - a f'(a)\bigr).
\]
La limite $f(a) - a f'(a)$ représente exactement \textbf{l'ordonnée à l'origine de la tangente $T_a$ à $\mathcal{C}_f$ au point d'abscisse $a$}. \hfill \textbf{[0{,}5 point]}

\medskip
\textbf{Méthode 2 : Approche par Développement Limité (Niveau MPSI)}

Puisque $f$ est dérivable en $a$, la formule de Taylor-Young à l'ordre 1 (ou définition infinitésimale de la dérivée) s'écrit au voisinage de $a$ :
\[
f(x) = f(a) + f'(a)(x - a) + o(x - a) \quad \text{quand } x \to a.
\]
Injectons ce développement limité dans le numérateur :
\[
x f(a) - a f(x) = x f(a) - a\Bigl(f(a) + f'(a)(x - a) + o(x - a)\Bigr) = (x - a)f(a) - a f'(a)(x - a) + o(x - a).
\]
Factorisons par $(x - a)$ :
\[
x f(a) - a f(x) = (x - a)\Bigl[ f(a) - a f'(a) \Bigr] + o(x - a).
\]
En divisant par $x - a \neq 0$ :
\[
\frac{x f(a) - a f(x)}{x - a} = f(a) - a f'(a) + o(1) \xrightarrow[x \to a]{} \mathbf{f(a) - a f'(a)}.
\] \hfill \textbf{[0{,}5 point]}
\end{bonusbox}

\end{document}
